\uuid{PozC}
\exo7id{3728}
\titre{exo7 3728}
\auteur{quercia}
\organisation{exo7}
\datecreate{2010-03-11}
\isIndication{false}
\isCorrection{true}
\chapitre{Espace euclidien, espace normé}
\sousChapitre{Forme quadratique}
\cpge{mpsi-mp2i.calcul-matriciel-systemes.e, pcsi.determinants.c, ptsi.determinants.a, mpsi-mp2i.determinants.d, pcsi.calcul-matriciel-systemes-lineaires.d, ptsi.calcul-matriciel-systemes-lineaires.d, bcpst1.matrices, mp-mpi.endomorphismes-euclidien.g, pc.complements-espaces-vectoriels.b, psi.complements-espaces-vectoriels.b, pt.complements-algebre-lineaire.b, bcpst2.valeurs-propres-vecteurs-propres.a}
\module{Algèbre}
\niveau{L2}
\difficulte{}

\contenu{
\texte{
Soit $A = (a_{i,j})\in \mathcal{M}_n(\R)$ à coefficients tous non nuls.
On note $A'$ la matrice de coefficient général $1/a_{i,j}$.
}
\begin{enumerate}
    \item \question{Trouver les matrices $A$ telles que $A$ et $A'$ sont symétriques définies positives
    (examiner les cas $n=1$, $n=2$, $n=n$).}
\reponse{Il n'y a pas de solution pour $n=2$ donc pas non plus pour $n>2$.}
    \item \question{Trouver les matrices $A$ telles que $A$ et $A'$ sont symétriques positives
    (examiner les cas $n=2$, $n=3$, $n=n$).}
\reponse{Pour $n=3$ on trouve $A={}^tCC$ où $C$ est une colonne sans zéros,
pour $n\ge 3$ on obtient le même résultat en considérant les blocs $3\times3$
centrés sur la diagonale.}
\end{enumerate}
}
